\section{Code}

Inline Code: \mintinline{python}{print("Hello World")}

\begin{listing}[ht]
	\begin{center}
		\begin{minted}{python}
print("Hello World")
        \end{minted}
	\end{center}
	\caption{Code Block}
	\label{fig:code-block}
\end{listing}

\begin{listing}[ht]
	\begin{center}
		\begin{minted}[
            linenos,
            % Font size (e.g. \small or \footnotesize if text is too big)
            fontsize=\normalsize,
            % Space between the lines
            baselinestretch=1.2,
            % Break long lines to a new line instead of overflowing
            breaklines,
            % Frame/Rules drawn around code block (single=rectangel, lines=top/bottom rule)
            frame=lines,
            % Frame/Rule color
            rulecolor=\color{codeframe},
            % Spacing between frame and conent
            framesep=2mm,
            % Left and right margin (before line numbers, after frame)
            xleftmargin=5mm,
            xrightmargin=5mm,
        ]{python}
# Make sure to not have any space before the code
def fibonacci(n):
    if n <= 1:
        return n
    return fibonacci(n-1) + fibonacci(n-2)
        \end{minted}
	\end{center}
	\caption{Code Block mit Zeilennummern und Abgrenzungen \texttt{[linenos, frame=lines]}}
\end{listing}

\begin{listing}[H]
	\begin{center}
		\inputminted[
			linenos,
			% Font size (e.g. \small or \footnotesize if text is too big)
			fontsize=\normalsize,
			% Space between the lines
			baselinestretch=1.2,
			% Break long lines to a new line instead of overflowing
			breaklines,
			% Frame/Rules drawn around code block (single=rectangel, lines=top/bottom rule)
			frame=lines,
			% Frame/Rule color
			rulecolor=\color{codeframe},
			% Spacing between frame and conent
			framesep=2mm,
			% Left and right margin (before line numbers, after frame)
			xleftmargin=5mm,
			xrightmargin=5mm,
			% Highlight specific lines
			highlightlines={14},
			% First line to show
			firstline=7,
		]{python}{code/server.py}

	\end{center}
	\caption{Code Block von Datei mit hervorgehobenen Zeilen}
\end{listing}

\begin{listing}[H]
	\begin{center}
		\inputminted[
			linenos,
			% Font size (e.g. \small or \footnotesize if text is too big)
			fontsize=\normalsize,
			% Space between the lines
			baselinestretch=1.2,
			% Break long lines to a new line instead of overflowing
			breaklines,
			% Frame/Rules drawn around code block (single=rectangel, lines=top/bottom rule)
			frame=lines,
			% Frame/Rule color
			rulecolor=\color{codeframe},
			% Spacing between frame and conent
			framesep=2mm,
			% Left and right margin (before line numbers, after frame)
			xleftmargin=5mm,
			xrightmargin=5mm,
			% Highlight specific lines
			highlightlines={3,4},
			% First line to show
			firstline=7,
			% First line number
			firstnumber=1,
		]{python}{code/client.py}

	\end{center}
	\caption{Code Block von Datei mit hervorgehobenen Zeilen}
\end{listing}
